% ===========================================================================
%  sviluppatore/06-dashboard/01-anatomia.tex
% ===========================================================================
\section{Anatomia di \pkgDashboard}\label{sez:anatomia-dashboard}

Trentasei file Python, \num{35} template Jinja, un centinaio di file
JavaScript e CSS. È l'oggetto più grande del progetto, e il modo più rapido
per orientarsi è guardare le dimensioni: dicono dove sta il peso.

\begin{center}
\begin{tabular}{@{}L{42mm}R{16mm}L{80mm}@{}}
\toprule
\textbf{File} & \textbf{KB} & \textbf{Di che cosa si occupa} \\
\midrule
\rigaTenue
\file{app.py}            & \num{226} & Il montaggio e \num{137} rotte. \\
\file{agent.py}          & \num{132} & L'assistente che scrive artefatti. \\
\rigaTenue
\file{community.py}      & \num{132} & Marketplace, contestazioni, account,
                                       messaggi, abbonamento. \\
\file{live\_runners.py}  & \num{90}  & I runner live: avvio, conto legato,
                                       kill-switch del conto, registro. \\
\rigaTenue
\file{web\_db.py}        & \num{90}  & L'accesso al Postgres del sito. \\
\file{agent\_tools\_sistema.py} & \num{80} & Gli strumenti con cui l'agente
                                       agisce sul sistema vivo. \\
\rigaTenue
\file{auth.py}           & \num{71}  & Accesso, registrazione, guardia
                                       globale, pagine HTML in linea. \\
\file{autopubblica.py}   & \num{49}  & Pubblicazione automatica degli
                                       artefatti, in un thread. \\
\rigaTenue
\file{genetic.py}        & \num{49}  & Strategie per programmazione genetica. \\
\file{visual\_debug.py}  & \num{45}  & Il debug visivo sui grafici. \\
\rigaTenue
\file{studio.py}         & \num{42}  & Editor, grafo a nodi, generazione. \\
\bottomrule
\end{tabular}
\end{center}

\subsection{Le quattro generazioni di codice, tutte vive}

Guardando come una rotta è registrata si capisce quando è stata scritta, e
quale stile seguire per estenderla.

\begin{center}
\begin{tabular}{@{}L{32mm}L{40mm}L{66mm}@{}}
\toprule
\textbf{Generazione} & \textbf{Dove} & \textbf{Come si riconosce} \\
\midrule
\rigaTenue
1 --- tutto in \file{app.py}
  & \cod{@app.route} dentro \simbolo{create\_app} & \num{137} rotte in un
    file di \num{226}~KB. \\
2 --- blueprint
  & \file{studio}, \file{community}, \file{agent}, \file{genetic}
  & \cod{Blueprint(...)}, nessun \cod{url\_prefix}. \\
\rigaTenue
3 --- estensioni v2
  & \file{ui\_ext\_*.py} & \simbolo{register(app, data\_layer, base\_path)},
    caricate da sole. \\
4 --- il guscio
  & \pkgApi{} + \pkgApplicazione & Non è nella dashboard: è l'altra
    applicazione (\vediSez{strangler}). \\
\bottomrule
\end{tabular}
\end{center}

\begin{pratica}{Quale stile usare oggi}
Per una rotta nuova che serve a una pagina: un modulo \file{ui\_ext\_*.py}.
Per una capacità nuova del sistema: un caso d'uso nel guscio, esposto
dall'API, chiamato dalla dashboard. In nessun caso si aggiunge una rotta a
\file{app.py}: è il file che tutte e tre le generazioni successive sono
servite a non far crescere.
\end{pratica}

\subsection{I template}

Trentadue pagine, tutte figlie di \file{base.html}, più tre frammenti
inclusi e \file{serve\_account.html}, che non è servito da nessuna rotta ma
dalla guardia di autenticazione.

Le pagine di accesso sono l'eccezione: \file{auth.py} non usa i template, ma
HTML in linea renderizzato con \simbolo{render\_template\_string}
(\cod{\_LOGIN\_HTML}, \cod{\_REGISTER\_HTML}, \cod{\_VERIFY\_HTML},
\cod{\_RESET\_HTML}, \cod{\_ATTESA\_HTML}). Chi cerca la pagina di login in
\file{templates/} non la trova.

\begin{nota}{Dove sta l'interfaccia del Live}
\file{live.html} non è servito da \percorsoHTTP{/live}: è servito da
\percorsoHTTP{/studio/runner-live}, cioè dal blueprint \cod{studio}. È il nome di
percorso più fuorviante del progetto (\vediSez{rotte}).
\end{nota}
